\documentclass[11pt,a4paper,oneside,openright]{book}
\usepackage[spanish,activeacute,es-noshorthands]{babel}
\usepackage[utf8]{inputenc}
\usepackage[T1]{fontenc}
\usepackage{lmodern}
\usepackage{amsmath, amssymb, amsthm}
\usepackage{graphicx}
\usepackage[dvipsnames]{xcolor}
\usepackage{tikz}
\usepackage{listings}
\usepackage{geometry}
\geometry{margin=2.5cm, headheight=14pt}
\usepackage{microtype}
\usepackage{booktabs}
\usepackage{longtable}
\usepackage{enumitem}
\usepackage{fancyhdr}
\usepackage{tcolorbox}
\tcbuselibrary{skins, breakable}
\usepackage{caption}
\captionsetup{font=small, labelfont=bf}
\usepackage{hyperref}
\hypersetup{
  colorlinks=true, linkcolor=NavyBlue, citecolor=ForestGreen, urlcolor=BrickRed,
  pdftitle={M11 --- CM5 + carrier: interaccion completa},
  pdfauthor={INTISAT}
}

\newtheorem{definicion}{Definicion}[chapter]
\newtheorem*{nota}{Nota}

\definecolor{codebg}{RGB}{248, 249, 250}
\definecolor{codekey}{RGB}{86, 156, 214}
\definecolor{codecomment}{RGB}{106, 153, 85}
\definecolor{codestr}{RGB}{206, 145, 120}

\lstdefinestyle{python}{
  language=Python,
  basicstyle=\ttfamily\footnotesize,
  keywordstyle=\color{codekey}\bfseries,
  commentstyle=\color{codecomment}\itshape,
  stringstyle=\color{codestr},
  numberstyle=\tiny\color{gray},
  numbers=left, numbersep=8pt,
  frame=single, rulecolor=\color{gray!30},
  framesep=4pt,
  backgroundcolor=\color{codebg},
  showstringspaces=false, breaklines=true, breakatwhitespace=true,
  tabsize=2, captionpos=b,
  literate={a}{{a}}1 {e}{{e}}1 {i}{{i}}1 {o}{{o}}1 {u}{{u}}1 {n}{{n}}1
}
\lstdefinestyle{shell}{
  basicstyle=\ttfamily\footnotesize,
  backgroundcolor=\color{codebg},
  frame=single, rulecolor=\color{gray!30},
  framesep=4pt, breaklines=true,
  showstringspaces=false,
  literate={a}{{a}}1 {e}{{e}}1 {i}{{i}}1 {o}{{o}}1 {u}{{u}}1 {n}{{n}}1
}
\lstset{style=python}

\newtcolorbox{intuibox}[1][]{
  colback=green!5, colframe=green!40!black,
  fonttitle=\bfseries, title=Intuicion,
  breakable, sharp corners=south, #1
}
\newtcolorbox{calculo}[1][]{
  colback=gray!4, colframe=gray!55,
  fonttitle=\bfseries, title=Calculo,
  breakable, sharp corners=south, #1
}
\newtcolorbox{payload}[1][]{
  colback=blue!4, colframe=blue!50!black,
  fonttitle=\bfseries, title=Conexion con INTISAT,
  breakable, sharp corners=south, #1
}
\newtcolorbox{cuidado}[1][]{
  colback=red!5, colframe=red!50!black,
  fonttitle=\bfseries, title=Cuidado,
  breakable, sharp corners=south, #1
}
\newtcolorbox{ejercicio}[1][]{
  colback=orange!5, colframe=orange!60!black,
  fonttitle=\bfseries, title=Ejercicios,
  breakable, sharp corners=south, #1
}

\pagestyle{fancy}
\fancyhf{}
\fancyhead[LE,RO]{\thepage}
\fancyhead[RE]{\nouppercase{\leftmark}}
\fancyhead[LO]{\nouppercase{\rightmark}}

\title{\Huge\bfseries M11 --- Compute Module 5\\[0.4em]
\Large Interaccion completa con el resto del payload INTISAT\\[0.6em]
\large Manual de la coleccion CubeSat (M11 / 11)}
\author{Coleccion CubeSat para el proyecto INTISAT}
\date{\today}

\begin{document}

\frontmatter
\maketitle

\chapter*{Prefacio}
\addcontentsline{toc}{chapter}{Prefacio}

Este manual responde una pregunta concreta: \emph{cuando reemplacemos el
Raspberry Pi 5 de desarrollo por el Compute Module 5 montado en un carrier
custom, exactamente como interactua con el resto del payload INTISAT?}

No es un manual de PCB design (eso esta en \texttt{M10\_CM5\_Carrier\_INTISAT.pdf}).
No es un manual de electronica basica (eso esta en
\texttt{Manual\_Electronica\_Practica.pdf}). No es un curso de Linux
(eso esta en \texttt{Raspberry\_Pi\_5\_Profundo.pdf}). Este manual cubre las
\textbf{interacciones entre el CM5 y todo lo que lo rodea}: alimentacion,
comunicacion con el OBC, control del sensor de imagen, control del OLED,
monitoreo de potencia y boot sequence. Lo justo que necesitas saber para
entender el sistema completo en funcionamiento.

\section*{Audiencia}

\begin{itemize}[leftmargin=2em]
  \item Ingenieros del payload que vienen del prototipo Pi 5 y quieren
        migrar al CM5.
  \item Cualquier miembro del equipo que pregunte "como anda eso".
  \item Operadores que tienen que diagnosticar problemas remotos.
  \item Lideres tecnicos que deben validar el diseño.
\end{itemize}

\section*{Prerequisitos}

\begin{itemize}[leftmargin=2em]
  \item \texttt{Manual\_Electronica\_Practica.pdf} (capitulos 1--12).
  \item \texttt{Raspberry\_Pi\_5\_Profundo.pdf} (capitulos 1--7).
  \item \texttt{M10\_CM5\_Carrier\_INTISAT.pdf} (al menos el cap. 1).
  \item Familiaridad con Python intermedio y Linux command line basico.
\end{itemize}

\section*{Como leer}

El manual son 12 capitulos mas 3 apendices. Total $\sim 35$ paginas.
Lectura completa: 1 dia laboral. Lectura selectiva: capitulos sueltos
segun necesidad.

\textbf{Cajas de color}:
\begin{itemize}[leftmargin=2em]
  \item \textbf{Verde} (Intuicion): analogias.
  \item \textbf{Gris} (Calculo): formulas y derivaciones.
  \item \textbf{Azul} (Conexion con INTISAT): como aplica a tu codigo.
  \item \textbf{Rojo} (Cuidado): errores comunes.
  \item \textbf{Naranja} (Ejercicios): preguntas para verificar.
\end{itemize}

\tableofcontents

\mainmatter

% =====================================================================
\chapter{Vision general del sandwich CM5 + carrier}
\label{ch:vision}

\section{Por que importa entender esto}

El Raspberry Pi 5 retail es un sistema \emph{cerrado}: vos enchufas el
USB-C y todo funciona. El \textbf{CM5 es un sistema abierto}: viene en
formato modulo, sin alimentacion propia, sin conectores externos. Es
\textbf{tu trabajo} diseñar la PCB que provea todo eso: el \emph{carrier
board}.

Entender la interaccion CM5 + carrier es entender como pasa de "computadora
en una caja" a "subsistema integrado en un satelite".

\section{La forma fisica del sandwich}

\begin{center}
\begin{tikzpicture}[scale=0.9]
% CM5
\draw[fill=blue!20, thick] (0,2) rectangle (6,3);
\node at (3, 2.5) {\textbf{Compute Module 5} ($55 \times 40$ mm)};
% Mezzanine
\draw[fill=gray!40] (0.3, 1.8) rectangle (2.8, 2.0);
\draw[fill=gray!40] (3.2, 1.8) rectangle (5.7, 2.0);
\node[font=\tiny] at (1.55, 1.9) {DF40C-100DS J1};
\node[font=\tiny] at (4.45, 1.9) {DF40C-100DS J2};
% Carrier
\draw[fill=green!10, thick] (-1, 0) rectangle (7, 1.8);
\node at (3, 0.9) {\textbf{Carrier board} ($\sim 90 \times 90$ mm)};
% Conectores carrier
\draw[fill=yellow!30] (0, -0.3) rectangle (1.2, 0);
\node[font=\tiny] at (0.6, -0.15) {PC-104 H1};
\draw[fill=yellow!30] (1.5, -0.3) rectangle (2.7, 0);
\node[font=\tiny] at (2.1, -0.15) {PC-104 H2};
\draw[fill=yellow!30] (3, -0.3) rectangle (4, 0);
\node[font=\tiny] at (3.5, -0.15) {FFC CSI-2};
\draw[fill=yellow!30] (4.3, -0.3) rectangle (5, 0);
\node[font=\tiny] at (4.65, -0.15) {SPI OLED};
\draw[fill=yellow!30] (5.3, -0.3) rectangle (5.9, 0);
\node[font=\tiny] at (5.6, -0.15) {Debug};
% Cotas
\draw[<->] (-1.3, 1.8) -- (-1.3, 3);
\node[rotate=90, font=\tiny] at (-1.5, 2.4) {1.5 mm stack};
\end{tikzpicture}
\end{center}

El CM5 se enchufa hacia abajo en los dos conectores hembra DF40 montados
en el carrier. Quedan a 1.5 mm de separacion (stack height nominal del
DF40C-100DS-0.4V).

\begin{intuibox}
Pensa el CM5 como un \emph{modulo SO-DIMM} de una notebook. Vos no diseñas
la SO-DIMM, la enchufas en la motherboard. La motherboard es el carrier.
Vos diseñas la motherboard.
\end{intuibox}

\section{Las 5 funciones del carrier}

\begin{enumerate}[leftmargin=2em]
  \item \textbf{Recibir alimentacion} del rail 5V\_PAYLOAD del INTISAT
        (via conector PC-104).
  \item \textbf{Regular esa alimentacion} y proteger al CM5 (LDO, fusible,
        diodo Schottky).
  \item \textbf{Conectar el CM5 a los buses externos}: I2C al OBC, CSI-2 al
        sensor, SPI al OLED.
  \item \textbf{Monitorear el consumo} con un sensor INA219.
  \item \textbf{Proveer debug y boot}: UART exposed, boton de boot, LED de
        status.
\end{enumerate}

\section{El stack PC-104 del INTISAT}

El INTISAT es un CubeSat 3U. Internamente tiene un stack de PCBs apiladas
con el estandar PC-104:

\begin{center}
\begin{tabular}{ll}
\toprule
\textbf{Posicion en el stack} & \textbf{Funcion} \\
\midrule
Top      & Payload (microscopia) --- tu carrier con CM5 \\
Middle 1 & ADCS (control de actitud) \\
Middle 2 & OBC central + COMMS \\
Bottom   & EPS (paneles, bateria, distribucion) \\
\bottomrule
\end{tabular}
\end{center}

Los 104 pines del PC-104 distribuyen todos los rails (5V, 3.3V, GND) y
los buses (I2C\_1, I2C\_2, CAN, UART) entre todas las cards del stack.

% =====================================================================
\chapter{Los 200 pines del mezzanine}
\label{ch:mezzanine}

\section{El conector DF40 de Hirose}

El CM5 expone toda su conectividad via 2 conectores hembra
\textbf{Hirose DF40C-100DS-0.4V}:

\begin{center}
\begin{tabular}{ll}
\toprule
\textbf{Parametro} & \textbf{Valor} \\
\midrule
Pines              & 100 por conector $\times$ 2 = 200 totales \\
Paso               & 0.4 mm \\
Stack height       & 1.5 mm (nominal) \\
Corriente max      & 0.5 A por pin \\
Voltaje max        & 30 V \\
Ciclos de mate     & 50 (no apto para conectores frecuentes) \\
Encapsulado        & SMD, soldar con horno por reflow \\
Precio             & USD 4 cada uno \\
\bottomrule
\end{tabular}
\end{center}

\subsection{Donde comprarlos}

\begin{itemize}[leftmargin=2em]
  \item \textbf{Mouser}: 855-DF40C-100DS-0.4V(81), USD 4.13.
  \item \textbf{Digikey}: H125185CT-ND, USD 4.07.
  \item \textbf{Arrow}: stock variable.
\end{itemize}

\begin{cuidado}
Si soldas el DF40 al reves (rotado 180°), el CM5 NO encaja. Marca la
orientacion en el silkscreen y verifica antes de soldar.
\end{cuidado}

\section{Categorias de pines}

Los 200 pines del CM5 se dividen funcionalmente:

\begin{center}
\begin{tabular}{p{4cm}p{2cm}p{8cm}}
\toprule
\textbf{Categoria} & \textbf{Cantidad} & \textbf{Que llevan} \\
\midrule
Power IN              & 4   & 5V de entrada principal \\
Power OUT             & 4   & 3V3 y 1V8 (que el CM5 te da para perifericos) \\
GND                   & 30  & tierra, multiples para baja impedancia \\
GPIO general          & 28  & los pines del header de la Pi 5 (BCM 0-27) \\
USB                   & 16  & 4 puertos USB (2 USB 3.0 + 2 USB 2.0) \\
Ethernet              & 8   & PHY signals (no incluye magnetics) \\
HDMI                  & 19  & video out (TMDS, I2C, hot plug) \\
DSI (display)         & 18  & 2 puertos para LCD/OLED grande \\
CSI (camera)          & 18  & 2 puertos camera (4 lanes c/u) \\
PCIe Gen 2            & 11  & 1 lane diferencial + clk + reset \\
SDIO                  & 8   & para SD externa (no usado en CM5 con eMMC) \\
RTC                   & 2   & VBAT + I2C del RTC \\
Control               & 6   & nRUN (reset), nPOWER\_BTN, EEPROM\_nDIS, LED, etc. \\
Reservado / NC        & 8   & no conectar \\
\bottomrule
\end{tabular}
\end{center}

\section{Cuales uso y cuales no para el INTISAT}

\begin{center}
\small
\begin{tabular}{p{3.5cm}p{1cm}p{1cm}p{7cm}}
\toprule
\textbf{Funcion} & \textbf{Uso?} & \textbf{Pines} & \textbf{Justificacion} \\
\midrule
5V IN               & Si   & 4   & alimentacion principal del CM5 \\
GND                 & Si   & todos & tierra de retorno \\
3V3 OUT             & Si   & 2-3 & alimentar TXS0108E, INA219, OLED \\
I2C1 (GPIO 2/3)     & Si   & 2   & primario, hacia OBC \\
SPI0 (GPIO 8-11)    & Si   & 4   & control del OLED SSD1351 \\
GPIO 18, 24         & Si   & 2   & DC y RST del OLED \\
CSI-2 (4 lanes)     & Si   & 18  & camara OV5647 \\
USB                 & No   & 16  & innecesario en orbita \\
HDMI                & No   & 19  & innecesario en orbita \\
Ethernet            & No   & 8   & innecesario en orbita \\
DSI                 & No   & 18  & no necesitamos display \\
PCIe                & No   & 11  & no necesitamos SSD NVMe \\
SDIO                & No   & 8   & usamos eMMC interna \\
nRUN (reset)        & Si   & 1   & para reset hardware desde el carrier \\
RTC VBAT            & Si   & 1   & opcional, mantener hora con bateria de boton \\
\bottomrule
\end{tabular}
\end{center}

\begin{payload}
Lo importante: \textbf{de los 200 pines, solo usas $\sim 50$}. Los otros 150
quedan no-conectados en tu carrier. Esto te da flexibilidad para diseñar
una PCB chica.
\end{payload}

\section{El pinout completo (extracto)}

Los siguientes son los pines mas usados para tu carrier. Pinout completo:
\url{https://datasheets.raspberrypi.com/cm5/cm5-datasheet.pdf}.

\begin{center}
\small
\begin{tabular}{llll}
\toprule
\textbf{Pin J1} & \textbf{Funcion} & \textbf{Pin J2} & \textbf{Funcion} \\
\midrule
1   & GND               & 1   & GND \\
3,5 & 5V\_IN            & 3,5 & 5V\_IN \\
7,9 & 3V3\_OUT          & 7,9 & 3V3\_OUT \\
22  & GPIO2 (I2C1 SDA)  & --- & --- \\
24  & GPIO3 (I2C1 SCL)  & --- & --- \\
26  & GPIO8 (SPI0 CE0)  & --- & --- \\
28  & GPIO9 (SPI0 MISO) & --- & --- \\
30  & GPIO10 (SPI0 MOSI)& --- & --- \\
32  & GPIO11 (SPI0 SCLK)& --- & --- \\
50  & GPIO18 (PWM/DC)   & --- & --- \\
54  & GPIO24            & --- & --- \\
--- & ---               & 71-89 & CSI-2 cam0 (4 lanes + clk) \\
99,100 & nRUN, GLOBAL\_EN & --- & --- \\
\bottomrule
\end{tabular}
\end{center}

\textbf{Importante}: estos son los pines tipicos. La asignacion exacta
varia con la version del firmware del CM5. \emph{Verificar siempre contra
el datasheet oficial}.

% =====================================================================
\chapter{Bloque de potencia detallado}
\label{ch:potencia}

\section{El camino completo de la energia}

Desde los paneles solares del INTISAT hasta los nucleos del CM5, la
energia atraviesa 7 etapas:

\begin{center}
\begin{tikzpicture}[scale=0.85, node distance=0.5cm, every node/.style={font=\small}]
\node[draw, fill=yellow!20] (sol) {Paneles solares};
\node[draw, fill=yellow!20, right=of sol] (mppt) {MPPT};
\node[draw, fill=yellow!20, right=of mppt] (bat) {Bateria Li-ion};
\node[draw, fill=blue!10, below=of bat] (rail) {Rail 5V\_PAYLOAD};
\node[draw, fill=red!10, left=of rail] (fuse) {Fusible PTC};
\node[draw, fill=red!10, left=of fuse] (diodo) {Diodo Schottky};
\node[draw, fill=red!10, left=of diodo] (ldo) {LDO 5V/6A};
\node[draw, fill=green!10, below=of ldo] (ina) {INA219};
\node[draw, fill=green!10, right=of ina] (cm5) {CM5 5V IN};
\node[draw, fill=green!10, right=of cm5] (pmic) {PMIC interna};

\draw[->] (sol) -- (mppt);
\draw[->] (mppt) -- (bat);
\draw[->] (bat) -- (rail);
\draw[->] (rail) -- (fuse);
\draw[->] (fuse) -- (diodo);
\draw[->] (diodo) -- (ldo);
\draw[->] (ldo) -- (ina);
\draw[->] (ina) -- (cm5);
\draw[->] (cm5) -- (pmic);
\end{tikzpicture}
\end{center}

Tu carrier toma el rail 5V\_PAYLOAD del PC-104 y, antes de pasar al CM5,
le aplica 3 elementos de proteccion + 1 de medicion.

\section{Etapa 1: Fusible PTC reseteable}

\begin{itemize}[leftmargin=2em]
  \item \textbf{Funcion}: si hay un corto en el CM5 o en algo aguas abajo,
        este fusible se "abre" (alta resistencia) y corta la corriente.
  \item \textbf{Tipo}: PTC (Positive Temperature Coefficient), reseteable.
        No se quema; se calienta y se enfria.
  \item \textbf{Eleccion}: \texttt{MF-MSMF050-2} de Bourns, 5 A hold,
        10 A trip. Cuesta USD 1.
  \item \textbf{Comportamiento}: a corriente nominal su resistencia es
        $\sim 0.02\,\Omega$ (negligible). En falla, $> 100\,\Omega$.
\end{itemize}

\begin{calculo}
Con 5 V y un corto que demanda 20 A, el PTC sube su resistencia a $0.25\,\Omega$:
\[
V_{\text{PTC}} = I \cdot R = 20 \cdot 0.25 = 5\;\text{V}
\]
Toda la tension cae sobre el PTC, la carga aguas abajo ve 0 V (apagada
hasta que el PTC se enfrie).
\end{calculo}

\section{Etapa 2: Diodo Schottky de proteccion polaridad inversa}

\begin{itemize}[leftmargin=2em]
  \item \textbf{Funcion}: si alguien conecta la fuente al reves, este
        diodo bloquea. Sin el, el CM5 se quema.
  \item \textbf{Tipo}: Schottky (bajo dropout, $\sim 0.3$ V vs $0.7$ V de
        Si normal).
  \item \textbf{Eleccion}: \texttt{1N5824} (3 A) o \texttt{SS54} (5 A).
        Cuesta USD 0.50.
  \item \textbf{Caida}: $V_F \approx 0.3$ V a 1 A. Disipacion: $0.3$ W.
\end{itemize}

\section{Etapa 3: LDO 5V o $\mu$Module}

Despues del diodo entran $\sim 4.7$ V (5 V - 0.3 V del diodo). El CM5
acepta $5.0$ V $\pm 5\%$, asi que necesitamos volver a $5.0$ V con poca
ripple.

\subsection{Opcion A: LDO}

\begin{itemize}[leftmargin=2em]
  \item \textbf{Chip}: \texttt{LT1086CT-5}, 1.5 A, dropout 1.3 V (no sirve
        a 4.7 V de entrada).
  \item \textbf{Mejor}: \texttt{TPS7A47}, ultra-low-noise, dropout
        $\sim 0.5$ V (apenas alcanza).
  \item \textbf{Problema}: pierde la diferencia entre $V_{in}$ y $V_{out}$
        como calor.
\end{itemize}

\subsection{Opcion B: $\mu$Module (recomendada)}

\begin{itemize}[leftmargin=2em]
  \item \textbf{Chip}: \texttt{LTM4631} de Analog Devices.
  \item \textbf{Tipo}: buck-boost (sube si entra bajo, baja si entra alto).
  \item \textbf{Entrada}: 2.7 V a 17 V.
  \item \textbf{Salida}: 5 V $@$ 6 A.
  \item \textbf{Eficiencia}: $94\%$ tipica.
  \item \textbf{Tamaño}: $9 \times 9 \times 4.92$ mm BGA.
  \item \textbf{Precio}: USD 25.
\end{itemize}

\begin{calculo}
Con LTM4631 a $94\%$ y CM5 consumiendo $4$ W:
\[
P_{\text{disipada}} = P_{\text{out}} \cdot \frac{1 - \eta}{\eta} = 4 \cdot \frac{0.06}{0.94} = 0.26\;\text{W}
\]
Con LDO equivalente, disiparia (4.7 - 5.0) $\cdot$ 0.8 A pero como no
puede subir voltaje, falla. Por eso el $\mu$Module gana.
\end{calculo}

\section{Etapa 4: INA219 monitor de potencia}

Ver \autoref{ch:ina219} para detalles. Resumen: mide voltaje, corriente y
potencia del rail 5V que entra al CM5, comunicada por I2C secundario.

\section{Etapa 5: Pin 5V\_IN del CM5}

Los 4 pines de 5V\_IN del CM5 (J1 pin 3, 5 y J2 pin 3, 5) tienen que
recibir la alimentacion regulada. Capacitor de bulk de $10\,\mu$F y
decoupling de $100$ nF en cada uno.

\section{Etapa 6: PMIC interna del CM5}

El CM5 tiene una PMIC (Power Management IC) que genera todos los rails
internos:

\begin{center}
\begin{tabular}{lll}
\toprule
\textbf{Rail} & \textbf{Voltaje} & \textbf{Para que} \\
\midrule
$V_{\text{CORE}}$  & 0.85--0.9 V & nucleos del CPU (DVFS variable) \\
$V_{\text{GPU}}$   & 0.9 V       & GPU VideoCore VII \\
$V_{\text{DDR}}$   & 1.1 V       & RAM LPDDR4X \\
$V_{\text{IO\_18}}$& 1.8 V       & buses internos del SoC \\
$V_{\text{IO\_33}}$& 3.3 V       & GPIO + perifericos externos \\
\bottomrule
\end{tabular}
\end{center}

Estos rails internos NO son accesibles desde el carrier (excepto el 3V3
que se expone via mezzanine pin 7, 9). Son uso exclusivo del CM5.

\section{Capacitores de decoupling}

\textbf{Regla universal}: cada chip con voltaje propio necesita un
capacitor de $100$ nF lo mas cerca posible.

En tu carrier:
\begin{itemize}[leftmargin=2em]
  \item $100$ nF en cada pin 5V\_IN del CM5 (4 unidades).
  \item $100$ nF en cada pin 3V3 que use perifericos (TXS, INA219, OLED).
  \item Bulk $10\,\mu$F al inicio del rail 5V (despues del LTM4631).
  \item Bulk $10\,\mu$F al inicio del rail 3V3 (lo que provee el CM5).
\end{itemize}

\begin{cuidado}
Sin decoupling, los chips fallan intermitentemente bajo carga rapida.
Es la regla numero uno de PCB design y se viola facilmente cuando uno
empieza. Verificar en el layout antes de mandar a fabricar.
\end{cuidado}

\section{Power-on sequence}

Cuando se aplica 5V al carrier:

\begin{enumerate}[leftmargin=2em]
  \item $t = 0$: rail 5V estable en el carrier.
  \item $t = 1$ ms: PMIC del CM5 genera $V_{\text{DDR}}$ (1.1 V).
  \item $t = 5$ ms: PMIC genera $V_{\text{CORE}}$ y $V_{\text{IO}}$.
  \item $t = 10$ ms: el SoC sale de reset.
  \item $t = 100$ ms: el bootloader EEPROM empieza.
  \item $t = 1$ s: bootloader carga \texttt{start.elf}.
  \item $t = 5$ s: kernel Linux empieza a cargar.
  \item $t = 30$ s: systemd termina de levantar servicios.
  \item $t = 31$ s: tu daemon esta listo, I2C slave en 0x42.
\end{enumerate}

\begin{ejercicio}
1) Calcular la potencia disipada por el LTM4631 cuando el CM5 esta en
PROCESSING (7.5 W). 2) Si el rail 5V\_PAYLOAD del INTISAT entrega
solo 4.5 V (caida en cables), verifica que el LTM4631 sigue dando
5.0 V regulados.
\end{ejercicio}

% =====================================================================
\chapter{Bus I2C con el OBC del INTISAT}
\label{ch:i2c}

\section{Concepto y nivel logico}

El OBC central del INTISAT y el CM5 hablan por I2C. Pero pueden estar
operando a distintos voltajes:

\begin{itemize}[leftmargin=2em]
  \item OBC: tipicamente 3.3 V o 5 V (depende del fabricante).
  \item CM5: 3.3 V en GPIO (no tolera 5 V directo).
\end{itemize}

Para conectarlos sin daño, se usa un \textbf{level shifter bidireccional}.

\section{TXS0108E: el level shifter elegido}

\begin{center}
\begin{tabular}{ll}
\toprule
\textbf{Parametro} & \textbf{Valor} \\
\midrule
Fabricante           & Texas Instruments \\
Canales              & 8 (usamos 2: SDA y SCL) \\
Voltaje lado A       & 1.2 V a 3.6 V \\
Voltaje lado B       & 1.65 V a 5.5 V \\
Direccion            & Auto (no necesita pin de control) \\
Velocidad max        & 110 Mbps \\
Encapsulado          & TSSOP-20 o WQFN-20 \\
Part number          & TXS0108EPWR \\
Precio               & USD 1 \\
\bottomrule
\end{tabular}
\end{center}

\subsection{Conexion en el carrier}

\begin{center}
\begin{tikzpicture}[scale=0.9, every node/.style={font=\small}]
\node[draw, fill=blue!10] (cm5) at (0, 1) {CM5 GPIO2/3 @ 3.3V};
\node[draw, fill=gray!10] (txs) at (4, 1) {TXS0108E};
\node[draw, fill=yellow!10] (obc) at (8, 1) {OBC SDA2/SCL2 @ 3.3V o 5V};
\draw[<->] (cm5) -- node[above] {SDA, SCL} (txs);
\draw[<->] (txs) -- node[above] {SDA, SCL} (obc);
\node[font=\tiny] at (2, 0.3) {pull-up 4.7k a 3.3V};
\node[font=\tiny] at (6, 0.3) {pull-up 4.7k a Vbus};
\end{tikzpicture}
\end{center}

\subsection{Configuracion del enable}

El TXS0108E tiene un pin \texttt{OE} (Output Enable). Ata a 3V3 con un
resistor de $10$ k$\Omega$ y un capacitor de $1\,\mu$F a GND para
soft-start. Sin este RC, el shifter puede dar glitches al power-on.

\section{Direccion I2C del payload}

El I2C slave del CM5 responde en \texttt{0x42} (decimal 66). Esta direccion
se eligio porque:

\begin{itemize}[leftmargin=2em]
  \item NO esta reservada por el estandar I2C.
  \item NO la usa ningun chip estandar (sensores comunes usan 0x40-0x4F,
        0x68-0x6F, etc.).
  \item Es facil de recordar.
\end{itemize}

\textbf{Verificacion}: ejecutar \texttt{i2cdetect -y 1} en el CM5 NO debe
mostrar 0x42 (es slave, no master). Pero en el OBC, su \texttt{i2cdetect}
si debe verlo.

\section{El BSC: I2C slave en el Pi}

El I2C estandar de Linux solo soporta master mode. Para hacer slave, el
Pi 5 / CM5 tiene un periferico especial llamado \textbf{BSC} (Broadcom
Serial Controller), accesible solo via \texttt{pigpio}.

\subsection{Codigo del slave (cubesat/i2c\_slave.py)}

\begin{lstlisting}
import pigpio
import time
from cubesat.commands import CMD_GET_STATUS, CMD_START_CAPTURE
from cubesat.commands import encode_status

I2C_ADDR = 0x42

def main():
    pi = pigpio.pi()
    if not pi.connected:
        print("ERROR: pigpiod not running")
        return

    # Configurar BSC en modo slave
    pi.bsc_i2c(I2C_ADDR)

    while True:
        # Lee bytes recibidos del master (OBC)
        status, byte_count, data = pi.bsc_i2c(I2C_ADDR)

        if byte_count > 0:
            cmd = data[0]
            payload = data[1:byte_count]

            # Despachar segun comando
            if cmd == CMD_GET_STATUS:
                response = encode_status(...)
                pi.bsc_i2c(I2C_ADDR, response)
            elif cmd == CMD_START_CAPTURE:
                # Encolar comando en /run/cubesat/commands/
                pass
            # ... mas comandos

        time.sleep(0.01)  # poll cada 10 ms

if __name__ == "__main__":
    main()
\end{lstlisting}

\section{Pull-up resistors}

I2C requiere pull-up resistors en SDA y SCL porque ambos son
\textbf{open-drain}: los chips solo pueden tirar la linea a 0 V; nadie
puede ponerla a $V_{cc}$.

\textbf{Valor tipico}: 4.7 k$\Omega$ a $V_{cc}$ (3.3 V en el lado A del
TXS, $V_{bus}$ en el lado B).

\begin{calculo}
Con 4.7 k$\Omega$ y 3.3 V:
\[
I_{\text{pullup}} = \frac{3.3}{4700} = 0.7\;\text{mA}
\]
A 100 kHz I2C, la capacitancia parasitica es $\sim 50$ pF. El tiempo de
subida tipico es $R \cdot C = 235$ ns, ok para 100 kHz (periodo $10\,\mu$s).
\end{calculo}

\section{Velocidad de operacion}

El I2C estandar opera a:

\begin{itemize}[leftmargin=2em]
  \item \textbf{100 kHz} (Standard mode) --- lo que usamos por defecto.
  \item \textbf{400 kHz} (Fast mode) --- si el OBC lo soporta y los cables
        son cortos.
  \item \textbf{1 MHz} (Fast mode plus) --- raro, requiere drivers
        especiales.
\end{itemize}

A 100 kHz: cada bit toma 10 $\mu$s. Un byte (9 bits con ACK) toma 90 $\mu$s.
Un thumbnail de 100 KB toma $\sim 10$ segundos a transmitir.

\section{Codigo del master en el OBC simulado}

Para probar desde tu laptop o desde otra Pi:

\begin{lstlisting}
import smbus2
import time

I2C_ADDR = 0x42
CMD_GET_STATUS = 0x01

bus = smbus2.SMBus(1)

# Enviar CMD_GET_STATUS
bus.write_byte(I2C_ADDR, CMD_GET_STATUS)
time.sleep(0.01)

# Leer respuesta (18 bytes: STATUS_OK + len + 16 bytes de status struct)
data = bus.read_i2c_block_data(I2C_ADDR, 0, 18)
print("Status:", data[0])
print("Length:", data[1])
print("Payload:", data[2:18])
\end{lstlisting}

\begin{payload}
Este codigo es el que vas a usar para validar el bring-up del carrier
antes de integrarlo al INTISAT. Lo corres desde otra Pi conectada a tu
carrier via SDA/SCL y verifica que el daemon responde.
\end{payload}

% =====================================================================
\chapter{Interfaz CSI-2 al sensor OV5647}
\label{ch:csi}

\section{Que es CSI-2}

\textbf{CSI-2} = Camera Serial Interface 2. Es el estandar de MIPI
Alliance para conectar camaras a procesadores. Lo usan todos los
smartphones y todas las camaras de Raspberry Pi.

\section{Caracteristicas electricas}

\begin{itemize}[leftmargin=2em]
  \item Hasta 4 lanes diferenciales de datos.
  \item 1 lane diferencial de clock.
  \item Voltaje: 200 mV pico a pico (D-PHY).
  \item Velocidad: 1.5 Gbps por lane.
  \item Impedancia: $90\,\Omega$ diferencial (con tolerancia $\pm 10\%$).
\end{itemize}

Para el OV5647 (2 lanes + 1 clock):
\begin{itemize}[leftmargin=2em]
  \item DAT0+, DAT0-
  \item DAT1+, DAT1-
  \item CLK+, CLK-
  \item Mas: I2C control (\texttt{CAM\_SDA}, \texttt{CAM\_SCL})
  \item Mas: poder ($V_{ANA}$, $V_{IO}$, $V_{DIG}$)
  \item Mas: nRESET, PWDN, LED\_CTRL
\end{itemize}

Total $\sim 15$ lineas.

\section{El cable FFC}

OV5647 se conecta al carrier por un cable plano flexible (FFC, Flat
Flexible Cable):

\begin{itemize}[leftmargin=2em]
  \item Estandar: 15-pin FFC, paso 1 mm.
  \item Longitud tipica: 15 cm.
  \item En el carrier: conector FFC vertical de 15 pines (\texttt{Molex 5025810-15}).
\end{itemize}

\begin{cuidado}
El conector FFC tiene polaridad. Los contactos pueden estar arriba o
abajo. La OV5647 v1.3 tiene los contactos abajo. Tu carrier debe tener
el conector con los contactos arriba para que el cable conecte correcto.
Si te equivocas, no daña pero no funciona.
\end{cuidado}

\section{Routing en el carrier}

Las lineas CSI-2 son criticas:

\begin{itemize}[leftmargin=2em]
  \item Cada par (DAT0+, DAT0-) debe rutearse \textbf{adyacente} y de la
        misma longitud.
  \item Match de longitud entre los dos pares: $\leq 1$ mm.
  \item Impedancia diferencial: $90\,\Omega$ (configurar en el stackup).
  \item Layer cercana al ground plane.
  \item Sin vias en el medio del par si es posible.
\end{itemize}

\textbf{Distancia maxima del CM5 al conector FFC en el carrier}: $\leq 100$ mm
(idealmente $\leq 50$ mm).

\section{Configuracion en Raspberry Pi OS}

En \texttt{/boot/firmware/config.txt}:

\begin{lstlisting}
# Habilitar camara
camera_auto_detect=1
dtoverlay=ov5647

# Opcional: I2C de la camara en bus distinto al primario
dtparam=i2c_vc=on
\end{lstlisting}

Despues del boot:

\begin{lstlisting}
libcamera-hello --list-cameras
# Debe mostrar: 0 : ov5647 [...]
\end{lstlisting}

\section{Codigo de captura (Python)}

Usando \texttt{Picamera2}:

\begin{lstlisting}
from picamera2 import Picamera2
import numpy as np

cam = Picamera2()

# Configuracion para microscopia (sin auto-correcciones)
config = cam.create_still_configuration(
    main={"size": (2592, 1944), "format": "RGB888"},
    raw={"size": (2592, 1944)}
)
cam.configure(config)

# Desactivar ISP para data cruda
cam.set_controls({
    "AwbEnable": False,           # auto white balance OFF
    "AeEnable": False,             # auto exposure OFF
    "ExposureTime": 50000,         # 50 ms
    "AnalogueGain": 1.0,           # sin amplificacion
    "ColourGains": (1.0, 1.0),     # ganancia R/B = 1
})
cam.start()

# Capturar
img = cam.capture_array("main")    # numpy array (H, W, 3)
print(f"Shape: {img.shape}, dtype: {img.dtype}")
\end{lstlisting}

\begin{payload}
En tu daemon \texttt{cubesat/capture.py}, esto se llama 25 veces por scan
FPM, una por cada angulo de iluminacion del OLED.
\end{payload}

% =====================================================================
\chapter{SPI al OLED SSD1351}
\label{ch:spi}

\section{SPI: 4 lineas + 2 de control}

\begin{itemize}[leftmargin=2em]
  \item \textbf{SCLK}: Serial Clock (GPIO11).
  \item \textbf{MOSI}: Master Out, Slave In (GPIO10). Datos del CM5 al OLED.
  \item \textbf{MISO}: Master In, Slave Out (GPIO9). \textit{No usado} en OLED.
  \item \textbf{CE0}: Chip Enable (GPIO8). El OLED escucha solo cuando CE0 baja.
  \item \textbf{DC}: Data/Command (GPIO18). 0 = comando, 1 = data.
  \item \textbf{RST}: Reset (GPIO24). Activo bajo.
\end{itemize}

\section{Conector en el carrier}

Header de 9 pines en el carrier (paso 2.54 mm) hacia el OLED:

\begin{center}
\begin{tabular}{cll}
\toprule
\textbf{Pin} & \textbf{Señal} & \textbf{Tension} \\
\midrule
1 & 3V3       & 3.3 V \\
2 & GND       & 0 V \\
3 & MOSI      & 3.3 V logic \\
4 & SCLK      & 3.3 V logic \\
5 & CE0       & 3.3 V logic \\
6 & DC        & 3.3 V logic \\
7 & RST       & 3.3 V logic \\
8 & V\_OLED   & 3.3 V de alimentacion del panel (puede requerir 12 V boost interno) \\
9 & GND       & 0 V \\
\bottomrule
\end{tabular}
\end{center}

\section{Configuracion en config.txt}

\begin{lstlisting}
# Habilitar SPI0
dtparam=spi=on
\end{lstlisting}

Despues:

\begin{lstlisting}
ls /dev/spidev*
# /dev/spidev0.0
\end{lstlisting}

\section{Codigo de control (Python)}

Usando \texttt{luma.oled}:

\begin{lstlisting}
from luma.core.interface.serial import spi
from luma.oled.device import ssd1351
from PIL import Image

# Setup SPI
serial = spi(port=0, device=0, gpio_DC=18, gpio_RST=24)
oled = ssd1351(serial, width=128, height=128)

# Crear imagen de iluminacion (un patron 5x5 con un pixel encendido)
img = Image.new("RGB", (128, 128), "black")
img.putpixel((64, 64), (255, 255, 255))    # pixel central encendido

# Mostrar
oled.display(img)
\end{lstlisting}

\begin{payload}
En FPM, se enciende uno solo de los 25 puntos posibles del 5x5 antes de
cada captura. Eso genera una "linterna" en diferente angulo. Tu codigo en
\texttt{cubesat/capture.py} ejecuta este loop 25 veces.
\end{payload}

\section{Velocidad de SPI}

El SPI0 puede operar de 0.5 MHz a 50 MHz. Para el OLED, 1 MHz es
suficiente y reduce EMI.

\section{Timing del Display}

El SSD1351 tarda $\sim 100$ ms en actualizar la pantalla completa
$128 \times 128$. Para FPM, encender un solo pixel toma menos: $\sim 10$ ms.

% =====================================================================
\chapter{Monitor INA219}
\label{ch:ina219}

\section{Que mide y por que importa}

El INA219 monitorea el rail 5V que alimenta al CM5. Provee:

\begin{itemize}[leftmargin=2em]
  \item \textbf{Voltaje del bus}: cuanto voltaje recibe el CM5.
  \item \textbf{Voltaje del shunt}: caida sobre el shunt de 0.1 $\Omega$.
  \item \textbf{Corriente}: calculada como $V_{\text{shunt}} / R_{\text{shunt}}$.
  \item \textbf{Potencia}: $V_{\text{bus}} \cdot I$.
\end{itemize}

\textbf{Aplicaciones}:

\begin{itemize}[leftmargin=2em]
  \item Telemetria: registrar consumo en cada modo (IDLE, CAPTURING, etc.).
  \item FDIR: detectar consumo anomalo (corto, regulador roto).
  \item Validacion: comparar consumo real vs estimado del datasheet.
\end{itemize}

\section{Especificaciones}

\begin{center}
\begin{tabular}{ll}
\toprule
\textbf{Parametro} & \textbf{Valor} \\
\midrule
Bus voltage max & 26 V \\
Shunt voltage max & $\pm 320$ mV \\
Corriente max & $\pm 3.2$ A (con shunt 0.1 $\Omega$) \\
Resolucion bus & 4 mV \\
Resolucion shunt & 10 $\mu$V \\
Resolucion corriente & 0.1 mA \\
Frecuencia conversion & hasta 1880 Hz \\
Comunicacion & I2C (4 address selectables) \\
Address default & 0x40 \\
Encapsulado & SOT-23-8 (chip) o modulo Adafruit \\
Precio (modulo) & USD 5 \\
\bottomrule
\end{tabular}
\end{center}

\section{Conexion al carrier}

\begin{center}
\begin{tikzpicture}[scale=0.9, every node/.style={font=\small}]
\node[draw, fill=yellow!20] (in) at (0, 1) {5V del LTM4631};
\node[draw, fill=red!10] (rs) at (3, 1) {Shunt 0.1$\Omega$};
\node[draw, fill=yellow!20] (out) at (6, 1) {5V al CM5};
\node[draw, fill=blue!10] (ina) at (3, -0.5) {INA219};
\draw[->] (in) -- (rs);
\draw[->] (rs) -- (out);
\draw[->] (in) |- (ina);
\draw[->] (out) |- (ina);
\node[font=\tiny] at (3, -1.5) {SDA, SCL al CM5 (I2C bus 2)};
\draw[->] (ina) -- (3, -2);
\end{tikzpicture}
\end{center}

\section{I2C bus separado}

El INA219 NO va en el mismo bus que el OBC. Razon: si el bus primario
se cae (por colision con master remoto, glitch electrico), el daemon
debe poder seguir leyendo telemetria.

Usar \textbf{I2C secundario} del CM5 (GPIO22/23 con dtoverlay
\texttt{i2c-gpio}):

\begin{lstlisting}
# en /boot/firmware/config.txt
dtoverlay=i2c-gpio,bus=2,i2c_gpio_sda=22,i2c_gpio_scl=23
\end{lstlisting}

\section{Codigo de lectura}

\begin{lstlisting}
from adafruit_ina219 import INA219
import board
import busio

i2c = busio.I2C(board.SCL, board.SDA)
ina = INA219(i2c)

while True:
    voltage = ina.bus_voltage + ina.shunt_voltage   # V real
    current = ina.current / 1000                     # convertir mA a A
    power = voltage * current                        # W

    print(f"V={voltage:.3f}V  I={current*1000:.1f}mA  P={power:.3f}W")
    time.sleep(0.5)
\end{lstlisting}

\section{Calibracion}

El INA219 tiene un registro de calibracion que define el factor de
escala para la corriente. Por defecto viene calibrado para shunt
$0.1\,\Omega$ y rango $\pm 3.2$ A. Si usas otro shunt, hay que reprogramar.

\begin{payload}
Tu \texttt{tools/power\_profiler.py} lee el INA219 cada 100 ms y guarda
el log en CSV durante un scan. Es lo que vamos a usar para validar el
presupuesto energetico real del payload.
\end{payload}

% =====================================================================
\chapter{Boot sequence}
\label{ch:boot}

\section{De power-on a daemon corriendo}

Cuando el INTISAT enciende el rail 5V\_PAYLOAD, el CM5 ejecuta esta
secuencia:

\begin{enumerate}[leftmargin=2em]
  \item \textbf{0--10 ms}: PMIC interna estabiliza rails.
  \item \textbf{10--100 ms}: bootloader EEPROM (ROM interno) inicializa
        el SoC.
  \item \textbf{100--500 ms}: bootloader carga el firmware del GPU desde
        la eMMC.
  \item \textbf{500 ms--2 s}: GPU inicializa LPDDR4X, lee \texttt{config.txt}
        y \texttt{cmdline.txt}.
  \item \textbf{2--5 s}: carga el kernel Linux desde la eMMC.
  \item \textbf{5--20 s}: kernel descomprime, descubre hardware, monta
        rootfs.
  \item \textbf{20--30 s}: systemd levanta servicios (pigpiod, networking,
        cubesat-*).
  \item \textbf{30--40 s}: tu daemon esta listo, I2C slave responde.
\end{enumerate}

\textbf{Total}: $\sim 30$ a $40$ segundos.

\section{Flashear la eMMC por primera vez}

Para programar el sistema operativo en la eMMC del CM5:

\begin{enumerate}[leftmargin=2em]
  \item Conectar el carrier por USB-C a una PC Linux.
  \item Mantener pulsado el boton de \texttt{nBOOT} del carrier mientras
        se aplica alimentacion.
  \item En la PC: \texttt{sudo rpiboot} (instalar previamente).
  \item La eMMC aparece como \texttt{/dev/sda} en la PC.
  \item Usar Raspberry Pi Imager para flashear Raspberry Pi OS Lite.
  \item Configurar usuario, SSH, WiFi (este ultimo se desactivara despues).
\end{enumerate}

\section{Debug por UART}

El CM5 expone una consola serial debug en GPIO14 (TX) y GPIO15 (RX).
Para usar:

\begin{enumerate}[leftmargin=2em]
  \item Conectar un adaptador USB-Serial (3.3 V) al header del carrier.
  \item Abrir \texttt{screen /dev/ttyUSB0 115200}.
  \item Ver los logs del boot en tiempo real.
\end{enumerate}

Esto es \textbf{crucial} en bring-up porque permite ver mensajes de error
antes que el sistema termine de arrancar.

% =====================================================================
\chapter{Flujo de datos durante un scan completo}
\label{ch:flujo}

\section{Timeline detallado}

Asumiendo que el OBC del INTISAT envia \texttt{CMD\_START\_CAPTURE} (modo
single):

\begin{center}
\small
\begin{tabular}{rp{12cm}}
\toprule
\textbf{Tiempo (ms)} & \textbf{Que pasa} \\
\midrule
0     & OBC inicia transmision I2C: \texttt{[START][0x42][0x02][0x00]} \\
2     & TXS0108E shifter traduce voltajes 5V $\to$ 3V3 \\
3     & CM5 GPIO2/3 reciben los bits \\
5     & pigpiod detecta byte recibido en BSC \\
10    & cubesat-i2c-slave service lee el comando \\
12    & escribe \texttt{cmd\_20260507\_101030.json} en \texttt{/run/cubesat/commands/} \\
15    & responde al OBC: \texttt{[STATUS\_OK]} \\
20    & cubesat-pipeline service detecta nuevo archivo (inotify) \\
30    & daemon cambia estado a CAPTURING \\
50    & enciende OLED en posicion (64, 64) via SPI \\
200   & OV5647 captura frame (50 ms exposicion + 150 ms processing) \\
250   & CM5 escribe \texttt{cap\_00.tiff} en \texttt{/var/cubesat/incoming/} \\
300   & atomic rename \texttt{cap\_00.tiff.tmp} $\to$ \texttt{cap\_00.tiff} \\
350   & daemon cambia estado a PROCESSING \\
400   & carga modelo StarDist en memoria \\
1500  & primera inferencia: segmenta el frame \\
2000  & calcula metricas (areas, perimeters) \\
2500  & escribe \texttt{summary.json}, \texttt{data.csv}, \texttt{overlay.png} \\
2800  & daemon cambia estado a EXPORTING \\
2900  & genera thumbnail 640x480 JPEG (\texttt{thumbnail.jpg}) \\
3000  & atomic write \texttt{downlink/thumbnail.jpg} \\
3050  & daemon vuelve a IDLE \\
3060  & telemetry service publica el evento en \texttt{/run/cubesat/telemetry.json} \\
\bottomrule
\end{tabular}
\end{center}

Total: $\sim 3$ segundos en modo single. Comparado con $\sim 60$ s en modo FPM.

\section{Datos generados}

Por cada scan single:

\begin{itemize}[leftmargin=2em]
  \item 1 frame raw TIFF: $\sim 5$ MB
  \item overlay.png: $\sim 500$ KB
  \item mask.png: $\sim 10$ KB
  \item data.csv: $\sim 5$ KB
  \item summary.json: $\sim 2$ KB
  \item thumbnail.jpg: $\sim 100$ KB
\end{itemize}

Total disco: $\sim 5.6$ MB por scan single.

\section{Downlink}

Cuando el OBC pide \texttt{CMD\_GET\_THUMBNAIL}:

\begin{enumerate}[leftmargin=2em]
  \item OBC: \texttt{write [0x42][0x06]}
  \item Daemon: lee thumbnail.jpg, lo trocea en chunks de 30 bytes.
  \item Daemon: responde con frames \texttt{[OK][30][more=1][chunk\_0..29]}
  \item OBC: lee 32 bytes por vez, acumula.
  \item Total: $\sim 3300$ frames para 100 KB.
  \item Tiempo total a 100 kHz: $\sim 10$ segundos.
\end{enumerate}

% =====================================================================
\chapter{Diferencias practicas Pi 5 retail vs CM5}
\label{ch:diferencias}

\section{Que cambia, que NO cambia}

\begin{center}
\begin{tabular}{p{4cm}p{5cm}p{5cm}}
\toprule
\textbf{Aspecto} & \textbf{Pi 5 retail} & \textbf{CM5 + carrier} \\
\midrule
Sistema operativo     & Raspberry Pi OS Lite Bookworm & idem \\
Kernel                & 6.6 + rpi & idem \\
Numeracion BCM        & GPIO 2, 3, etc. & idem \\
Codigo del daemon     & funciona & funciona (sin cambios) \\
Comandos shell        & funcionan & funcionan \\
Conectores fisicos    & USB-C + 4 USB-A + HDMI + Ethernet + header & via mezzanine al carrier \\
Almacenamiento        & SD card & eMMC soldada \\
Como entra la energia & USB-C exterior & rail 5V del PC-104 \\
WiFi/Ethernet         & si, con conectores & desactivados en config \\
Boot                  & power-on con SD montada & power-on con eMMC \\
Flashear primera vez  & SD a PC, escribir, insertar & USB-C a PC, rpiboot, flashear \\
Recovery              & quitar SD, escribir nueva & rpiboot otra vez \\
Tamaño                & 85$\times$56 mm & 55$\times$40 mm (mas tu carrier) \\
Peso                  & 46 g & 17 g + tu carrier \\
\bottomrule
\end{tabular}
\end{center}

\section{Lo unico que cambia en software}

\begin{enumerate}[leftmargin=2em]
  \item \texttt{config.txt}: agregar \texttt{dtoverlay=disable-wifi},
        \texttt{dtoverlay=disable-bt}, configurar HDMI off.
  \item \texttt{cubesat/install.sh}: detectar eMMC en lugar de SD.
  \item Documentacion: actualizar guia de flasheo (rpiboot en lugar de
        Raspberry Pi Imager directo).
\end{enumerate}

% =====================================================================
\chapter{Procedimiento de bring-up}
\label{ch:bringup}

Cuando llega el primer carrier de fabrica:

\section{Step 1: inspeccion visual (15 min)}

\begin{itemize}[leftmargin=2em]
  \item Microscopio o lupa 10x.
  \item Verificar todas las soldaduras: sin solder bridges, sin joints frios.
  \item Verificar componentes en sus footprints correctos.
  \item Verificar polaridad de capacitores polarizados (tantalio).
  \item Verificar orientacion de diodos (catodo marca).
  \item Verificar el silkscreen de los DF40 (matching del CM5).
\end{itemize}

\section{Step 2: tests electricos sin CM5 (30 min)}

Sin el CM5 montado, con una fuente externa:

\begin{enumerate}[leftmargin=2em]
  \item Multimetro en modo continuidad: verificar GND continuo en todos
        los pines GND del carrier.
  \item Multimetro en modo continuidad: verificar que NO hay corto entre
        5V y GND, entre 3V3 y GND.
  \item Alimentar con fuente $5$ V, limite $100$ mA. Verificar el rail
        5V despues del LTM4631 = 5.00 $\pm$ 0.05 V.
  \item Verificar que el INA219 responde:
        \texttt{sudo i2cdetect -y 2}, debe mostrar 0x40.
  \item Subir el limite a $1$ A, verificar que el rail no cae.
\end{enumerate}

\section{Step 3: montar el CM5 (15 min)}

\begin{itemize}[leftmargin=2em]
  \item Verificar orientacion: la flecha del CM5 alineada con la marca
        del carrier.
  \item Presionar uniformemente desde el centro hacia los bordes.
  \item Verificar stack height: 1.5 mm $\pm$ 0.2.
  \item Atornillar con M3 si el carrier tiene agujeros de fijacion.
\end{itemize}

\section{Step 4: primer boot via UART (30 min)}

\begin{enumerate}[leftmargin=2em]
  \item Conectar USB-Serial al header de debug del carrier.
  \item Abrir \texttt{screen /dev/ttyUSB0 115200}.
  \item Aplicar alimentacion 5 V.
  \item Ver los logs en tiempo real.
  \item Verificar que termina con prompt \texttt{login:}.
\end{enumerate}

\section{Step 5: validacion de buses (20 min)}

Desde SSH al CM5 (despues de configurar):

\begin{lstlisting}
# I2C primario (debe ver TXS0108E si tiene un sensor del lado externo)
sudo i2cdetect -y 1

# I2C secundario (debe ver INA219 en 0x40)
sudo i2cdetect -y 2

# SPI (lista de dispositivos)
ls /dev/spidev*

# Camara
libcamera-hello --list-cameras

# Lectura INA219
python -c "
from adafruit_ina219 import INA219
import board, busio
i2c = busio.I2C(board.SCL, board.SDA)
ina = INA219(i2c)
print(f'V={ina.bus_voltage:.2f} V, I={ina.current:.1f} mA')
"
\end{lstlisting}

\section{Step 6: instalar el daemon y validar (30 min)}

\begin{lstlisting}
cd /opt
sudo git clone https://github.com/JaminYC/CubeSat-EdgeAI-Payload.git
cd CubeSat-EdgeAI-Payload
sudo bash cubesat/install.sh

# Verificar servicios
systemctl status cubesat-i2c-slave
systemctl status cubesat-pipeline
systemctl status cubesat-telemetry

# Probar comando I2C desde otra Pi
# (desde la otra Pi)
i2cset -y 1 0x42 0x01
i2cget -y 1 0x42 0
\end{lstlisting}

\section{Step 7: test funcional end-to-end (60 min)}

Desde otra Pi (simulando el OBC):

\begin{enumerate}[leftmargin=2em]
  \item \texttt{CMD\_GET\_STATUS}: verificar respuesta valida.
  \item \texttt{CMD\_START\_CAPTURE} (modo single): verificar que se crea
        archivo en \texttt{/var/cubesat/incoming/}.
  \item Esperar el procesado.
  \item \texttt{CMD\_GET\_THUMBNAIL}: descargar y verificar que es un
        JPEG valido.
  \item Repetir 10 veces para validar repetibilidad.
\end{enumerate}

\textbf{Si los 7 steps pasan, el carrier esta listo para integrar al stack
PC-104 del INTISAT.}

% =====================================================================
\chapter{Troubleshooting comun}
\label{ch:troubleshooting}

\section{El CM5 no bootea}

\subsection{Sintoma: LED rojo encendido continuo, sin LED verde}

\begin{itemize}[leftmargin=2em]
  \item \textbf{Causa probable}: alimentacion insuficiente.
  \item \textbf{Verificar}: voltaje en pin 5V\_IN del CM5 con multimetro.
        Debe ser 5.00 $\pm$ 0.05 V.
  \item \textbf{Si el voltaje cae al cargar}: el LTM4631 no esta dando
        suficiente corriente. Verificar el inductor, los capacitores.
\end{itemize}

\subsection{Sintoma: LED verde parpadea N veces}

Codigos de error documentados en
\url{https://github.com/raspberrypi/firmware/issues}.

\begin{itemize}[leftmargin=2em]
  \item 4 parpadeos: no encuentra \texttt{start.elf} en la eMMC.
  \item 7 parpadeos: kernel image corrupta.
  \item 8 parpadeos: error de SDRAM.
\end{itemize}

\section{El I2C no responde}

\subsection{Sintoma: \texttt{i2cdetect} no muestra 0x42}

Recordar: el CM5 es slave, no aparece en su propio \texttt{i2cdetect}.
Probar desde otra Pi (master).

\subsection{Sintoma: el master ve 0x42 pero los comandos no responden}

\begin{itemize}[leftmargin=2em]
  \item Verificar que \texttt{pigpiod} esta corriendo:
        \texttt{systemctl status pigpiod}.
  \item Verificar que \texttt{cubesat-i2c-slave.service} esta running.
  \item Ver logs: \texttt{journalctl -u cubesat-i2c-slave -f}.
\end{itemize}

\section{La camara no se detecta}

\subsection{Sintoma: \texttt{libcamera-hello --list-cameras} no la muestra}

\begin{itemize}[leftmargin=2em]
  \item Verificar conexion del cable FFC: contactos hacia arriba en el
        conector del carrier.
  \item Verificar \texttt{config.txt}: \texttt{dtoverlay=ov5647}.
  \item Verificar la I2C de la camara:
        \texttt{sudo i2cdetect -y 10} (bus camera), debe ver 0x36 (OV5647).
\end{itemize}

\section{El OLED no muestra nada}

\begin{itemize}[leftmargin=2em]
  \item Verificar alimentacion 3V3 del OLED (multimetro).
  \item Verificar que el RST esta en alto (no en reset permanente).
  \item Probar con \texttt{luma.oled} ejemplo basico.
\end{itemize}

\section{El INA219 lee 0}

\begin{itemize}[leftmargin=2em]
  \item Verificar shunt 0.1 $\Omega$ presente y sin corto.
  \item Verificar calibracion del INA219 (registro 0x05).
  \item Verificar polaridad: IN+ del lado de la fuente, IN- del lado del CM5.
\end{itemize}

\section{El consumo es mas alto del esperado}

\begin{itemize}[leftmargin=2em]
  \item ¿WiFi desactivado? Verificar \texttt{rfkill list}.
  \item ¿BT desactivado? Verificar \texttt{hciconfig}.
  \item ¿HDMI off? Verificar \texttt{vcgencmd display\_power 0}.
  \item ¿Throttling? Verificar \texttt{vcgencmd get\_throttled}, no debe ser 0x0.
\end{itemize}

% =====================================================================
\appendix
\chapter{Apendice A --- Pinout completo del mezzanine}

Pines mas usados del CM5 para el INTISAT:

\begin{center}
\small
\begin{longtable}{lllll}
\toprule
\textbf{Pin J} & \textbf{Numero} & \textbf{Funcion} & \textbf{Tipo} & \textbf{Notas} \\
\midrule
\endhead
J1 & 1, 100      & GND            & power & ref electrica \\
J1 & 3, 5        & 5V\_IN         & power & alimentacion principal \\
J1 & 7, 9        & 3V3\_OUT       & power & 3.3 V del CM5 \\
J1 & 22          & GPIO2 (SDA1)   & I2C   & SDA bus 1 \\
J1 & 24          & GPIO3 (SCL1)   & I2C   & SCL bus 1 \\
J1 & 26          & GPIO8          & SPI   & CE0 \\
J1 & 28          & GPIO9          & SPI   & MISO \\
J1 & 30          & GPIO10         & SPI   & MOSI \\
J1 & 32          & GPIO11         & SPI   & SCLK \\
J1 & 50          & GPIO18         & GPIO  & DC OLED \\
J1 & 54          & GPIO24         & GPIO  & RST OLED \\
J1 & 99          & nRUN           & ctrl  & reset hardware (active low) \\
J1 & 100         & GLOBAL\_EN     & ctrl  & habilitacion global \\
J2 & 71          & CAM\_DAT0\_P   & CSI-2 & lane 0 positivo \\
J2 & 72          & CAM\_DAT0\_N   & CSI-2 & lane 0 negativo \\
J2 & 73          & CAM\_DAT1\_P   & CSI-2 & lane 1 positivo \\
J2 & 74          & CAM\_DAT1\_N   & CSI-2 & lane 1 negativo \\
J2 & 75          & CAM\_CLK\_P    & CSI-2 & clock positivo \\
J2 & 76          & CAM\_CLK\_N    & CSI-2 & clock negativo \\
J2 & 80          & CAM\_GPIO0     & I2C   & SDA camera \\
J2 & 82          & CAM\_GPIO1     & I2C   & SCL camera \\
\bottomrule
\end{longtable}
\end{center}

\chapter{Apendice B --- BOM del carrier (estimado)}

\begin{center}
\begin{longtable}{p{4cm}p{4cm}p{2.5cm}p{2cm}}
\toprule
\textbf{Componente} & \textbf{Part number} & \textbf{Proveedor} & \textbf{USD} \\
\midrule
\endhead
CM5 (4 GB/32 GB)                  & SC1230                 & RPi Trading & 60 \\
DF40 mezzanine (2x)                & DF40C-100DS-0.4V       & Hirose      & 8 \\
$\mu$Module 5V/6A                  & LTM4631                & Analog Dev  & 25 \\
Diodo Schottky 5A                  & SS54                   & Vishay      & 0.50 \\
Fusible PTC 5A reseteable          & MF-MSMF050             & Bourns      & 1 \\
INA219 modulo                      & INA219AIDR             & TI          & 3 \\
TXS0108E level shifter             & TXS0108EPWR            & TI          & 1 \\
Cap. ceramico 100 nF (50x)         & GRM188R71H104KA93D     & Murata      & 2.50 \\
Cap. tantalio 10 uF (10x)          & T491A106K016AT         & Kemet       & 3 \\
Conector FFC 15-pin                & 5025810-09             & Molex       & 1 \\
Conector header 7-pin OLED         & estandar 2.54 mm       & varios      & 0.50 \\
Boton SMD tactile                  & 1825910-6              & TE          & 0.50 \\
LED bicolor SMD                    & APBA3025               & Kingbright  & 0.50 \\
Resistor 5.6 ohm 1\%               & varios                 & varios      & 0.05 \\
Resistor 4.7 k$\Omega$ 1\% (10x)    & varios                 & varios      & 0.10 \\
PCB 4-layer (10 unidades)          & JLCPCB                 & ---         & 50 \\
\bottomrule
\end{longtable}
\end{center}

\textbf{Total para 1 unidad}: $\sim$ USD 156.

\chapter{Apendice C --- Checklist pre-integracion al PC-104}

\begin{itemize}[leftmargin=2em]
  \item[$\square$] Inspeccion visual: sin solder bridges, todos los chips en footprint correcto.
  \item[$\square$] Continuidad: GND continuo, sin cortos a rails de alimentacion.
  \item[$\square$] Test sin CM5: el LTM4631 da 5.00 $\pm$ 0.05 V a vacio y con $1$ A de carga.
  \item[$\square$] INA219 detectado en I2C bus 2 (0x40).
  \item[$\square$] CM5 montado correctamente con stack height 1.5 mm.
  \item[$\square$] Boot por UART completa hasta prompt de login.
  \item[$\square$] eMMC flasheada con Raspberry Pi OS Lite Bookworm.
  \item[$\square$] Configuracion aplicada: WiFi/BT/HDMI desactivados.
  \item[$\square$] Daemon \texttt{cubesat-pipeline} corriendo.
  \item[$\square$] I2C slave responde en 0x42 desde otra Pi (master).
  \item[$\square$] Camara OV5647 detectada por \texttt{libcamera-hello}.
  \item[$\square$] OLED SSD1351 responde a comandos via \texttt{luma.oled}.
  \item[$\square$] Test funcional: 10 scans en modo single sin error.
  \item[$\square$] OTA: subir un commit, verificar update + rollback.
  \item[$\square$] Watchdog systemd: simular cuelgue, verificar reset.
  \item[$\square$] Burn-in: 168 h corriendo sin reboot.
  \item[$\square$] Termico: T del SoC $< 70$ °C bajo carga sostenida.
  \item[$\square$] Documentacion: bring-up log firmado y archivado.
\end{itemize}

\backmatter

\chapter*{Bibliografia}
\addcontentsline{toc}{chapter}{Bibliografia}

\begin{enumerate}[leftmargin=2em]
\item Raspberry Pi Foundation. \emph{Compute Module 5 Datasheet}. \\
  \url{https://datasheets.raspberrypi.com/cm5/cm5-datasheet.pdf}, 2024.
\item Raspberry Pi Foundation. \emph{CM5 IO Board Datasheet}. \\
  \url{https://datasheets.raspberrypi.com/cm5/cm5-ioboard-datasheet.pdf}, 2024.
\item Hirose Electric. \emph{DF40 Series Catalog}. \\
  \url{https://www.hirose.com/en/product/series/DF40}.
\item Analog Devices. \emph{LTM4631 Datasheet}. \\
  \url{https://www.analog.com/en/products/ltm4631.html}.
\item Texas Instruments. \emph{TXS0108E Datasheet}. \\
  \url{https://www.ti.com/product/TXS0108E}.
\item Texas Instruments. \emph{INA219 Datasheet}. \\
  \url{https://www.ti.com/product/INA219}.
\item OmniVision. \emph{OV5647 Datasheet}, 2009.
\item Solomon Systech. \emph{SSD1351 Datasheet}, 2014.
\item MIPI Alliance. \emph{D-PHY Specification v2.5}, 2022.
\item Manual M0 --- CubeSat 101.
\item Manual M2 --- Sistemas de Potencia EPS.
\item Manual M5 --- Flight Software.
\item Manual M10 --- CM5 + Carrier (diseño del PCB).
\end{enumerate}

\end{document}
